% Original analytic teaching diagram, 2026-11-12.
% Splitting the binary digits d_1, d_2, ... of one uniform point into disjoint groups:
% group k uses the positions 2^{k-1}(2j-1), j = 1,2,...; group k builds U_k = 0.(digits of group k) in binary.
\definecolor{VB}{HTML}{7188A3}\definecolor{VBF}{HTML}{E5EAF0}
\definecolor{VG}{HTML}{829C8B}\definecolor{VGF}{HTML}{E8EEE9}
\definecolor{VY}{HTML}{B4A16F}\definecolor{VYF}{HTML}{F1ECDD}
\definecolor{VR}{HTML}{A7787D}\definecolor{VRF}{HTML}{F0E5E6}
\definecolor{VP}{HTML}{9588A5}\definecolor{VPF}{HTML}{ECE8F0}
\definecolor{VI}{HTML}{3B4252}\definecolor{VM}{HTML}{6B7280}\definecolor{VL}{HTML}{4C4D4F}
\begin{tikzpicture}[x=1mm,y=1mm,>=Stealth,font=\figurefont\fontsize{8}{10}\selectfont,
lab/.style={text=VM,font=\figurefont\fontsize{7.5}{9}\selectfont}]
\def\cw{7.6}
% group of each position 1..12
\foreach \p/\fc/\dc in {1/VBF/VB,2/VGF/VG,3/VBF/VB,4/VYF/VY,5/VBF/VB,6/VGF/VG,7/VBF/VB,8/VPF/VP,9/VBF/VB,10/VGF/VG,11/VBF/VB,12/VYF/VY} {
  \fill[\fc] ({(\p-1)*\cw},20) rectangle ({\p*\cw},26);
  \draw[\dc,line width=.7pt] ({(\p-1)*\cw},20) rectangle ({\p*\cw},26);
  \node[text=VI] at ({(\p-.5)*\cw},23) {$d_{\p}$};
}
\node[anchor=west,text=VM] at ({12*\cw+1},23) {$\cdots$};
\node[anchor=west,text=VB] at (0,14) {$U_1=0.d_1d_3d_5d_7\cdots$};
\node[anchor=west,text=VG] at (48,14) {$U_2=0.d_2d_6d_{10}d_{14}\cdots$};
\node[anchor=west,text=VY] at (0,6) {$U_3=0.d_4d_{12}d_{20}d_{28}\cdots$};
\node[anchor=west,text=VP] at (48,6) {$U_4=0.d_8d_{24}d_{40}d_{56}\cdots$};
\end{tikzpicture}
